% --- Archon physics formalization source begin ---
% archon:physics
% archon:covers PhyXMiniProblems/problem_phyx_mini_0859.lean
% archon:source-report reports/phyx_mini/problem_phyx_mini_0859.source.json

\chapter{Physics problem phyx\_mini\_0859}
\label{ch:PhyXMiniProblems_problem_phyx_mini_0859}

\paragraph{Problem source.}
\#\# Question

What is the electric potential at the point indicated with the dot?

\#\# Answer choices

- A: A: -1100V

- B: B: 1600V

- C: C: -800V

- D: D: -1600V

\#\# Figure caption (auxiliary; use the image as primary evidence)

The image shows an equilateral triangle with charges placed at each vertex. 

- At the top vertex, there is a red circle with a positive charge of \textbackslash{}(+1.0 \textbackslash{}, \textbackslash{}text\{nC\}\textbackslash{}).
- At the bottom left vertex, there is a cyan circle with a negative charge of \textbackslash{}(-2.0 \textbackslash{}, \textbackslash{}text\{nC\}\textbackslash{}).
- At the bottom right vertex, there is another cyan circle with a negative charge of \textbackslash{}(-2.0 \textbackslash{}, \textbackslash{}text\{nC\}\textbackslash{}).

All sides of the triangle are equal, and each side is labeled as \textbackslash{}(3.0 \textbackslash{}, \textbackslash{}text\{cm\}\textbackslash{}).

A small black dot is located at the center of the triangle, indicating the centroid or a point of interest within the triangle. 

Dashed lines connect the charges to indicate the sides of the triangle.

\#\# Dataset metadata

Electrostatics; Physical Model Grounding Reasoning, Spatial Relation Reasoning

\paragraph{Recorded answer/context.}
D: D: -1600V

\paragraph{Figure/image path.}
/root/proposal\_for\_physic/science-mango/phyx\_mini\_run/phyx\_data/test\_image/859.png

\paragraph{Formalization target.}
create a compiling Lean file with sorry bodies at `PhyXMiniProblems/problem\_phyx\_mini\_0859.lean`. The Lean declarations must preserve the physical quantities, dimensions or dimensional roles, figure labels, governing-law hypotheses, and final relation expressed by this problem.
Use Mathlib/Physlib names found through LeanExplore where available. If a physics API is missing, introduce faithful local abstractions rather than scalar placeholder aliases.

\begin{theorem}[Physics formalization target]
\label{thm:physics:phyx_mini_0859:target}
\lean{PhyXMiniProblems.ProblemPhyXMini0859.problem_phyx_mini_0859}
\uses{def:physics:phyx-mini-0859:phyxminiproblems-problemphyxmini0859-equilateralpointchargesetup, def:physics:phyx-mini-0859:phyxminiproblems-problemphyxmini0859-potentialatobservationinvolts, def:physics:phyx-mini-0859:phyxminiproblems-problemphyxmini0859-matchesprimaryfigure859, def:physics:phyx-mini-0859:phyxminiproblems-problemphyxmini0859-usesvacuumcoulombconstant, def:physics:phyx-mini-0859:phyxminiproblems-problemphyxmini0859-hasphysicaltriangleparameters, def:physics:phyx-mini-0859:phyxminiproblems-problemphyxmini0859-satisfiespointchargepotentialandsuperpositionlaws, def:physics:phyx-mini-0859:phyxminiproblems-problemphyxmini0859-displayedpotentialinvolts, def:physics:phyx-mini-0859:phyxminiproblems-problemphyxmini0859-roundstonearesthundredvolts, def:physics:phyx-mini-0859:phyxminiproblems-problemphyxmini0859-isuniquenearestdisplayedpotential}
The assigned autoformalize agent should translate this physics problem into Lean declarations in the covered file, with theorem and lemma proof bodies written as `by sorry`.
\end{theorem}
\begin{proof}
This is an autoformalization task, not a proof task. Produce statements that can later be proved by the physics prover without weakening the physical model.
\end{proof}

% --- Archon declaration topology begin ---
\section{Declaration topology}
\begin{definition}[electric Potential Dimension]
  \label{def:physics:phyx-mini-0859:phyxminiproblems-problemphyxmini0859-electricpotentialdimension}
  \lean{PhyXMiniProblems.ProblemPhyXMini0859.electricPotentialDimension}
  The physical dimension M L² T⁻² C⁻¹ of electric potential
\end{definition}
\begin{proof}
  This is the declared model definition of electric Potential Dimension.
\end{proof}
\begin{definition}[Length Quantity]
  \label{def:physics:phyx-mini-0859:phyxminiproblems-problemphyxmini0859-lengthquantity}
  \lean{PhyXMiniProblems.ProblemPhyXMini0859.LengthQuantity}
  A nonnegative, unit-independent physical length
\end{definition}
\begin{proof}
  This is the declared model definition of Length Quantity.
\end{proof}
\begin{definition}[Signed Charge Quantity]
  \label{def:physics:phyx-mini-0859:phyxminiproblems-problemphyxmini0859-signedchargequantity}
  \lean{PhyXMiniProblems.ProblemPhyXMini0859.SignedChargeQuantity}
  A signed, unit-independent physical electric charge
\end{definition}
\begin{proof}
  This is the declared model definition of Signed Charge Quantity.
\end{proof}
\begin{definition}[Electric Potential Quantity]
  \label{def:physics:phyx-mini-0859:phyxminiproblems-problemphyxmini0859-electricpotentialquantity}
  \lean{PhyXMiniProblems.ProblemPhyXMini0859.ElectricPotentialQuantity}
  \uses{def:physics:phyx-mini-0859:phyxminiproblems-problemphyxmini0859-electricpotentialdimension}
  A signed, unit-independent physical electric potential
\end{definition}
\begin{proof}
  This is the declared model definition of Electric Potential Quantity.
\end{proof}
\begin{definition}[length In Meters]
  \label{def:physics:phyx-mini-0859:phyxminiproblems-problemphyxmini0859-lengthinmeters}
  \lean{PhyXMiniProblems.ProblemPhyXMini0859.lengthInMeters}
  \uses{def:physics:phyx-mini-0859:phyxminiproblems-problemphyxmini0859-lengthquantity}
  Coherent-SI metre readout of a physical length
\end{definition}
\begin{proof}
  This is the declared model definition of length In Meters.
\end{proof}
\begin{definition}[length In Centimeters]
  \label{def:physics:phyx-mini-0859:phyxminiproblems-problemphyxmini0859-lengthincentimeters}
  \lean{PhyXMiniProblems.ProblemPhyXMini0859.lengthInCentimeters}
  \uses{def:physics:phyx-mini-0859:phyxminiproblems-problemphyxmini0859-lengthquantity, def:physics:phyx-mini-0859:phyxminiproblems-problemphyxmini0859-lengthinmeters}
  Centimetre readout used by the three printed side labels
\end{definition}
\begin{proof}
  This is the declared model definition of length In Centimeters.
\end{proof}
\begin{definition}[charge In Coulombs]
  \label{def:physics:phyx-mini-0859:phyxminiproblems-problemphyxmini0859-chargeincoulombs}
  \lean{PhyXMiniProblems.ProblemPhyXMini0859.chargeInCoulombs}
  \uses{def:physics:phyx-mini-0859:phyxminiproblems-problemphyxmini0859-signedchargequantity}
  Coherent-SI coulomb readout of a signed physical charge
\end{definition}
\begin{proof}
  This is the declared model definition of charge In Coulombs.
\end{proof}
\begin{definition}[charge In Nanocoulombs]
  \label{def:physics:phyx-mini-0859:phyxminiproblems-problemphyxmini0859-chargeinnanocoulombs}
  \lean{PhyXMiniProblems.ProblemPhyXMini0859.chargeInNanocoulombs}
  \uses{def:physics:phyx-mini-0859:phyxminiproblems-problemphyxmini0859-signedchargequantity, def:physics:phyx-mini-0859:phyxminiproblems-problemphyxmini0859-chargeincoulombs}
  Nanocoulomb readout used by the three printed charge labels
\end{definition}
\begin{proof}
  This is the declared model definition of charge In Nanocoulombs.
\end{proof}
\begin{definition}[electric Potential In Volts]
  \label{def:physics:phyx-mini-0859:phyxminiproblems-problemphyxmini0859-electricpotentialinvolts}
  \lean{PhyXMiniProblems.ProblemPhyXMini0859.electricPotentialInVolts}
  \uses{def:physics:phyx-mini-0859:phyxminiproblems-problemphyxmini0859-electricpotentialquantity}
  Coherent-SI volt readout of a signed physical electric potential
\end{definition}
\begin{proof}
  This is the declared model definition of electric Potential In Volts.
\end{proof}
\begin{definition}[Source Charge]
  \label{def:physics:phyx-mini-0859:phyxminiproblems-problemphyxmini0859-sourcecharge}
  \lean{PhyXMiniProblems.ProblemPhyXMini0859.SourceCharge}
  The three point charges, named by their locations in the image
\end{definition}
\begin{proof}
  This is the declared model definition of Source Charge.
\end{proof}
\begin{definition}[source Vertex Index]
  \label{def:physics:phyx-mini-0859:phyxminiproblems-problemphyxmini0859-sourcevertexindex}
  \lean{PhyXMiniProblems.ProblemPhyXMini0859.sourceVertexIndex}
  \uses{def:physics:phyx-mini-0859:phyxminiproblems-problemphyxmini0859-sourcecharge}
  Vertex index used for each named source in Mathlib's affine triangle
\end{definition}
\begin{proof}
  This is the declared model definition of source Vertex Index.
\end{proof}
\begin{definition}[Triangle Side]
  \label{def:physics:phyx-mini-0859:phyxminiproblems-problemphyxmini0859-triangleside}
  \lean{PhyXMiniProblems.ProblemPhyXMini0859.TriangleSide}
  The three dashed sides of the equilateral triangle
\end{definition}
\begin{proof}
  This is the declared model definition of Triangle Side.
\end{proof}
\begin{definition}[side Endpoints]
  \label{def:physics:phyx-mini-0859:phyxminiproblems-problemphyxmini0859-sideendpoints}
  \lean{PhyXMiniProblems.ProblemPhyXMini0859.sideEndpoints}
  \uses{def:physics:phyx-mini-0859:phyxminiproblems-problemphyxmini0859-sourcecharge, def:physics:phyx-mini-0859:phyxminiproblems-problemphyxmini0859-triangleside}
  The pair of charged vertices joined by each named side
\end{definition}
\begin{proof}
  This is the declared model definition of side Endpoints.
\end{proof}
\begin{definition}[Figure Charge Color]
  \label{def:physics:phyx-mini-0859:phyxminiproblems-problemphyxmini0859-figurechargecolor}
  \lean{PhyXMiniProblems.ProblemPhyXMini0859.FigureChargeColor}
  The red and cyan fills used for positive and negative charges
\end{definition}
\begin{proof}
  This is the declared model definition of Figure Charge Color.
\end{proof}
\begin{definition}[Figure Charge Sign]
  \label{def:physics:phyx-mini-0859:phyxminiproblems-problemphyxmini0859-figurechargesign}
  \lean{PhyXMiniProblems.ProblemPhyXMini0859.FigureChargeSign}
  The sign glyph drawn inside each coloured charge circle
\end{definition}
\begin{proof}
  This is the declared model definition of Figure Charge Sign.
\end{proof}
\begin{definition}[expected Charge Color]
  \label{def:physics:phyx-mini-0859:phyxminiproblems-problemphyxmini0859-expectedchargecolor}
  \lean{PhyXMiniProblems.ProblemPhyXMini0859.expectedChargeColor}
  \uses{def:physics:phyx-mini-0859:phyxminiproblems-problemphyxmini0859-sourcecharge, def:physics:phyx-mini-0859:phyxminiproblems-problemphyxmini0859-figurechargecolor}
  Expected source-circle colour read from the primary image
\end{definition}
\begin{proof}
  This is the declared model definition of expected Charge Color.
\end{proof}
\begin{definition}[expected Charge Sign]
  \label{def:physics:phyx-mini-0859:phyxminiproblems-problemphyxmini0859-expectedchargesign}
  \lean{PhyXMiniProblems.ProblemPhyXMini0859.expectedChargeSign}
  \uses{def:physics:phyx-mini-0859:phyxminiproblems-problemphyxmini0859-sourcecharge, def:physics:phyx-mini-0859:phyxminiproblems-problemphyxmini0859-figurechargesign}
  Expected sign glyph read from the primary image
\end{definition}
\begin{proof}
  This is the declared model definition of expected Charge Sign.
\end{proof}
\begin{definition}[expected Charge In Nanocoulombs]
  \label{def:physics:phyx-mini-0859:phyxminiproblems-problemphyxmini0859-expectedchargeinnanocoulombs}
  \lean{PhyXMiniProblems.ProblemPhyXMini0859.expectedChargeInNanocoulombs}
  \uses{def:physics:phyx-mini-0859:phyxminiproblems-problemphyxmini0859-sourcecharge}
  Signed nanocoulomb number printed beside each source circle
\end{definition}
\begin{proof}
  This is the declared model definition of expected Charge In Nanocoulombs.
\end{proof}
\begin{definition}[Equilateral Charge Figure]
  \label{def:physics:phyx-mini-0859:phyxminiproblems-problemphyxmini0859-equilateralchargefigure}
  \lean{PhyXMiniProblems.ProblemPhyXMini0859.EquilateralChargeFigure}
  \uses{def:physics:phyx-mini-0859:phyxminiproblems-problemphyxmini0859-sourcecharge, def:physics:phyx-mini-0859:phyxminiproblems-problemphyxmini0859-triangleside, def:physics:phyx-mini-0859:phyxminiproblems-problemphyxmini0859-figurechargecolor, def:physics:phyx-mini-0859:phyxminiproblems-problemphyxmini0859-figurechargesign}
  Literal presentation data transcribed from image 859.png
\end{definition}
\begin{proof}
  This is the declared model definition of Equilateral Charge Figure.
\end{proof}
\begin{definition}[Equilateral Point Charge Setup]
  \label{def:physics:phyx-mini-0859:phyxminiproblems-problemphyxmini0859-equilateralpointchargesetup}
  \lean{PhyXMiniProblems.ProblemPhyXMini0859.EquilateralPointChargeSetup}
  \uses{def:physics:phyx-mini-0859:phyxminiproblems-problemphyxmini0859-lengthquantity, def:physics:phyx-mini-0859:phyxminiproblems-problemphyxmini0859-signedchargequantity, def:physics:phyx-mini-0859:phyxminiproblems-problemphyxmini0859-electricpotentialquantity, def:physics:phyx-mini-0859:phyxminiproblems-problemphyxmini0859-sourcecharge, def:physics:phyx-mini-0859:phyxminiproblems-problemphyxmini0859-equilateralchargefigure}
  The independent electrostatic system. sourcePotential and totalPotential are physical scalar fields on the plane; their relationship is imposed only by the governing laws below. Neither field is defined from an answer choice
\end{definition}
\begin{proof}
  This is the declared model definition of Equilateral Point Charge Setup.
\end{proof}
\begin{definition}[source Position]
  \label{def:physics:phyx-mini-0859:phyxminiproblems-problemphyxmini0859-sourceposition}
  \lean{PhyXMiniProblems.ProblemPhyXMini0859.sourcePosition}
  \uses{def:physics:phyx-mini-0859:phyxminiproblems-problemphyxmini0859-sourcecharge, def:physics:phyx-mini-0859:phyxminiproblems-problemphyxmini0859-sourcevertexindex, def:physics:phyx-mini-0859:phyxminiproblems-problemphyxmini0859-equilateralpointchargesetup}
  Spatial position of a named source charge
\end{definition}
\begin{proof}
  This is the declared model definition of source Position.
\end{proof}
\begin{definition}[side Length In Meters]
  \label{def:physics:phyx-mini-0859:phyxminiproblems-problemphyxmini0859-sidelengthinmeters}
  \lean{PhyXMiniProblems.ProblemPhyXMini0859.sideLengthInMeters}
  \uses{def:physics:phyx-mini-0859:phyxminiproblems-problemphyxmini0859-triangleside, def:physics:phyx-mini-0859:phyxminiproblems-problemphyxmini0859-sideendpoints, def:physics:phyx-mini-0859:phyxminiproblems-problemphyxmini0859-equilateralpointchargesetup, def:physics:phyx-mini-0859:phyxminiproblems-problemphyxmini0859-sourceposition}
  Metre-coordinate distance between the endpoints of a named side
\end{definition}
\begin{proof}
  This is the declared model definition of side Length In Meters.
\end{proof}
\begin{definition}[source Distance To Observation In Meters]
  \label{def:physics:phyx-mini-0859:phyxminiproblems-problemphyxmini0859-sourcedistancetoobservationinmeters}
  \lean{PhyXMiniProblems.ProblemPhyXMini0859.sourceDistanceToObservationInMeters}
  \uses{def:physics:phyx-mini-0859:phyxminiproblems-problemphyxmini0859-sourcecharge, def:physics:phyx-mini-0859:phyxminiproblems-problemphyxmini0859-equilateralpointchargesetup, def:physics:phyx-mini-0859:phyxminiproblems-problemphyxmini0859-sourceposition}
  Metre-coordinate distance from a source charge to the observation dot
\end{definition}
\begin{proof}
  This is the declared model definition of source Distance To Observation In Meters.
\end{proof}
\begin{definition}[potential At Observation In Volts]
  \label{def:physics:phyx-mini-0859:phyxminiproblems-problemphyxmini0859-potentialatobservationinvolts}
  \lean{PhyXMiniProblems.ProblemPhyXMini0859.potentialAtObservationInVolts}
  \uses{def:physics:phyx-mini-0859:phyxminiproblems-problemphyxmini0859-electricpotentialinvolts, def:physics:phyx-mini-0859:phyxminiproblems-problemphyxmini0859-equilateralpointchargesetup}
  Volt readout of the resultant potential at the black dot
\end{definition}
\begin{proof}
  This is the declared model definition of potential At Observation In Volts.
\end{proof}
\begin{definition}[Matches Primary Figure859]
  \label{def:physics:phyx-mini-0859:phyxminiproblems-problemphyxmini0859-matchesprimaryfigure859}
  \lean{PhyXMiniProblems.ProblemPhyXMini0859.MatchesPrimaryFigure859}
  \uses{def:physics:phyx-mini-0859:phyxminiproblems-problemphyxmini0859-lengthinmeters, def:physics:phyx-mini-0859:phyxminiproblems-problemphyxmini0859-lengthincentimeters, def:physics:phyx-mini-0859:phyxminiproblems-problemphyxmini0859-chargeinnanocoulombs, def:physics:phyx-mini-0859:phyxminiproblems-problemphyxmini0859-expectedchargecolor, def:physics:phyx-mini-0859:phyxminiproblems-problemphyxmini0859-expectedchargesign, def:physics:phyx-mini-0859:phyxminiproblems-problemphyxmini0859-expectedchargeinnanocoulombs, def:physics:phyx-mini-0859:phyxminiproblems-problemphyxmini0859-equilateralpointchargesetup, def:physics:phyx-mini-0859:phyxminiproblems-problemphyxmini0859-sidelengthinmeters}
  All numerical and qualitative information read from the primary image. The physical observation point is identified with Mathlib's triangle centroid. No electric-potential value or answer choice occurs in this predicate
\end{definition}
\begin{proof}
  This is the declared model definition of Matches Primary Figure859.
\end{proof}
\begin{definition}[Uses Vacuum Coulomb Constant]
  \label{def:physics:phyx-mini-0859:phyxminiproblems-problemphyxmini0859-usesvacuumcoulombconstant}
  \lean{PhyXMiniProblems.ProblemPhyXMini0859.UsesVacuumCoulombConstant}
  \uses{def:physics:phyx-mini-0859:phyxminiproblems-problemphyxmini0859-equilateralpointchargesetup}
  Standard vacuum calibration of Physlib's Coulomb constant, in SI units
\end{definition}
\begin{proof}
  This is the declared model definition of Uses Vacuum Coulomb Constant.
\end{proof}
\begin{definition}[Has Physical Triangle Parameters]
  \label{def:physics:phyx-mini-0859:phyxminiproblems-problemphyxmini0859-hasphysicaltriangleparameters}
  \lean{PhyXMiniProblems.ProblemPhyXMini0859.HasPhysicalTriangleParameters}
  \uses{def:physics:phyx-mini-0859:phyxminiproblems-problemphyxmini0859-lengthinmeters, def:physics:phyx-mini-0859:phyxminiproblems-problemphyxmini0859-equilateralpointchargesetup, def:physics:phyx-mini-0859:phyxminiproblems-problemphyxmini0859-sourceposition}
  Positivity and noncoincidence conditions for the depicted setup
\end{definition}
\begin{proof}
  This is the declared model definition of Has Physical Triangle Parameters.
\end{proof}
\begin{definition}[Satisfies Point Charge Potential And Superposition Laws]
  \label{def:physics:phyx-mini-0859:phyxminiproblems-problemphyxmini0859-satisfiespointchargepotentialandsuperpositionlaws}
  \lean{PhyXMiniProblems.ProblemPhyXMini0859.SatisfiesPointChargePotentialAndSuperpositionLaws}
  \uses{def:physics:phyx-mini-0859:phyxminiproblems-problemphyxmini0859-chargeincoulombs, def:physics:phyx-mini-0859:phyxminiproblems-problemphyxmini0859-electricpotentialinvolts, def:physics:phyx-mini-0859:phyxminiproblems-problemphyxmini0859-sourcecharge, def:physics:phyx-mini-0859:phyxminiproblems-problemphyxmini0859-equilateralpointchargesetup, def:physics:phyx-mini-0859:phyxminiproblems-problemphyxmini0859-sourceposition}
  For a point charge q, the electrostatic potential at a point a distance r away is k q / r. Potentials superpose as signed scalars. The first field states the point-charge law at every nonsource point, and the second states superposition at every point; neither mentions the requested numerical value
\end{definition}
\begin{proof}
  This is the declared model definition of Satisfies Point Charge Potential And Superposition Laws.
\end{proof}
\begin{lemma}[source Distance To Centroid eq side div sqrt Three]
  \label{lem:physics:phyx-mini-0859:phyxminiproblems-problemphyxmini0859-sourcedistancetocentroid-eq-side-div-sqrtthree}
  \lean{PhyXMiniProblems.ProblemPhyXMini0859.sourceDistanceToCentroid_eq_side_div_sqrtThree}
  \uses{def:physics:phyx-mini-0859:phyxminiproblems-problemphyxmini0859-lengthinmeters, def:physics:phyx-mini-0859:phyxminiproblems-problemphyxmini0859-equilateralpointchargesetup, def:physics:phyx-mini-0859:phyxminiproblems-problemphyxmini0859-sourcedistancetoobservationinmeters, def:physics:phyx-mini-0859:phyxminiproblems-problemphyxmini0859-matchesprimaryfigure859, def:physics:phyx-mini-0859:phyxminiproblems-problemphyxmini0859-hasphysicaltriangleparameters}
  In an equilateral triangle, the centroid is also the circumcenter, and its distance from every vertex is the side length divided by √3
\end{lemma}
\begin{proof}
  The conclusion follows from the governing relations and figure readouts named in the statement of source Distance To Centroid eq side div sqrt Three.
\end{proof}
\begin{lemma}[potential At Observation eq point Charge Sum]
  \label{lem:physics:phyx-mini-0859:phyxminiproblems-problemphyxmini0859-potentialatobservation-eq-pointchargesum}
  \lean{PhyXMiniProblems.ProblemPhyXMini0859.potentialAtObservation_eq_pointChargeSum}
  \uses{def:physics:phyx-mini-0859:phyxminiproblems-problemphyxmini0859-chargeincoulombs, def:physics:phyx-mini-0859:phyxminiproblems-problemphyxmini0859-sourcecharge, def:physics:phyx-mini-0859:phyxminiproblems-problemphyxmini0859-equilateralpointchargesetup, def:physics:phyx-mini-0859:phyxminiproblems-problemphyxmini0859-sourcedistancetoobservationinmeters, def:physics:phyx-mini-0859:phyxminiproblems-problemphyxmini0859-potentialatobservationinvolts, def:physics:phyx-mini-0859:phyxminiproblems-problemphyxmini0859-hasphysicaltriangleparameters, def:physics:phyx-mini-0859:phyxminiproblems-problemphyxmini0859-satisfiespointchargepotentialandsuperpositionlaws}
  The inverse-distance law and scalar superposition express the potential at the dot as the sum of the three source contributions
\end{lemma}
\begin{proof}
  The conclusion follows from the governing relations and figure readouts named in the statement of potential At Observation eq point Charge Sum.
\end{proof}
\begin{lemma}[potential At Observation numerical Bounds]
  \label{lem:physics:phyx-mini-0859:phyxminiproblems-problemphyxmini0859-potentialatobservation-numericalbounds}
  \lean{PhyXMiniProblems.ProblemPhyXMini0859.potentialAtObservation_numericalBounds}
  \uses{def:physics:phyx-mini-0859:phyxminiproblems-problemphyxmini0859-equilateralpointchargesetup, def:physics:phyx-mini-0859:phyxminiproblems-problemphyxmini0859-potentialatobservationinvolts, def:physics:phyx-mini-0859:phyxminiproblems-problemphyxmini0859-matchesprimaryfigure859, def:physics:phyx-mini-0859:phyxminiproblems-problemphyxmini0859-usesvacuumcoulombconstant, def:physics:phyx-mini-0859:phyxminiproblems-problemphyxmini0859-hasphysicaltriangleparameters, def:physics:phyx-mini-0859:phyxminiproblems-problemphyxmini0859-satisfiespointchargepotentialandsuperpositionlaws}
  With the image charges, side length, and vacuum calibration, the actual potential lies strictly between -1650 V and -1550 V
\end{lemma}
\begin{proof}
  The conclusion follows from the governing relations and figure readouts named in the statement of potential At Observation numerical Bounds.
\end{proof}
\begin{definition}[Answer Choice]
  \label{def:physics:phyx-mini-0859:phyxminiproblems-problemphyxmini0859-answerchoice}
  \lean{PhyXMiniProblems.ProblemPhyXMini0859.AnswerChoice}
  Labels of the four electric-potential choices printed in the question
\end{definition}
\begin{proof}
  This is the declared model definition of Answer Choice.
\end{proof}
\begin{definition}[displayed Potential In Volts]
  \label{def:physics:phyx-mini-0859:phyxminiproblems-problemphyxmini0859-displayedpotentialinvolts}
  \lean{PhyXMiniProblems.ProblemPhyXMini0859.displayedPotentialInVolts}
  \uses{def:physics:phyx-mini-0859:phyxminiproblems-problemphyxmini0859-answerchoice}
  Electric potential in volts printed beside each answer label
\end{definition}
\begin{proof}
  This is the declared model definition of displayed Potential In Volts.
\end{proof}
\begin{definition}[Rounds To Nearest Hundred Volts]
  \label{def:physics:phyx-mini-0859:phyxminiproblems-problemphyxmini0859-roundstonearesthundredvolts}
  \lean{PhyXMiniProblems.ProblemPhyXMini0859.RoundsToNearestHundredVolts}
  A computed voltage rounds to a displayed value at the nearest 100 V
\end{definition}
\begin{proof}
  This is the declared model definition of Rounds To Nearest Hundred Volts.
\end{proof}
\begin{definition}[Is Unique Nearest Displayed Potential]
  \label{def:physics:phyx-mini-0859:phyxminiproblems-problemphyxmini0859-isuniquenearestdisplayedpotential}
  \lean{PhyXMiniProblems.ProblemPhyXMini0859.IsUniqueNearestDisplayedPotential}
  \uses{def:physics:phyx-mini-0859:phyxminiproblems-problemphyxmini0859-answerchoice, def:physics:phyx-mini-0859:phyxminiproblems-problemphyxmini0859-displayedpotentialinvolts}
  A displayed choice is strictly nearer to the computed voltage than every alternative
\end{definition}
\begin{proof}
  This is the declared model definition of Is Unique Nearest Displayed Potential.
\end{proof}
% --- Archon declaration topology end ---
% --- Archon physics formalization source end ---
